\setcounter{exercisegroup}{1}
\exercisegroup

\begin{enumerate}
\def\labelenumi{\arabic{enumi})}
\item
  A subset A of a vector space E, is starshaped relative to 0 if for all \(x \in \mathbf{A}\) and every \(\lambda\) such that \(0 \leqslant \lambda < 1\) , the point \(\lambda x\) belongs to A. Let A be starshaped and such that, for each \(x \in \mathbf{A}\) , there exists \(\mu > 1\) such that \(\mu x \in \mathbf{A}\) . Show that if, for every pair of points \(x, y\) of A we have \(\frac{1}{2}(x + y) \in \mathbf{A}\) , then A is convex. Give an example of a non-convex starshaped set A such that \(\frac{1}{2}(\mathbf{A} + \mathbf{A}) \subset \mathbf{A}\) .
\item
  Let A be a convex subset of an affine space E and B a set containing A. Show that, amongst the convex sets that both contain A and are contained in B there exists at least one maximal set; give an example where there are several distinct maximal sets.
\item
  Let A and B be two disjoint convex sets in a vector space E. Show that there exist two disjoint convex sets C, D in E such that \(A \subset C\) , \(B \subset D\) and \(C \cup D = E\) . (Apply th. 2 of S, III, § 2.4 to the set of pairs of disjoint convex sets (M, N) such that \(A \subset M\) and \(B \subset N\) and express the fact that M and N do not meet by the relation \(0 \notin M - N\) . To show that \(C \cup D = E\) , obtain a contradiction supposing that \(x_{0} \notin C \cup D\) ; if \(C'\) (resp. \(D'\) ) is the convex envelope of \(C \cup \{x_{0}\}\) (resp. \(D \cup \{x_{0}\}\) ), show that it is impossible that both \(C' \cap D \neq \varnothing\) and \(C \cap D' \neq \varnothing\) .)
\item
  Let C be a convex cone with vertex 0 in a vector space E; if \((x_{i})_{1\leqslant i\leqslant n}\) is a finite family of points of C such that \(\sum_{i = 1}^{n}\lambda_{i}x_{i} = 0\) for a family of numbers \(\lambda_{i} > 0\) , then C contains the vector subspace of E generated by the \(x_{i}\) .
\item
  Suppose that the vector space E has an enumerably infinite basis \((e_{n})_{n\in\mathbb{N}}\) . Let C be the set of points \(x = \sum_{n} \xi_{n} e_{n}\) such that for the largest index n for which \(\xi_{n} \neq 0\) , we have \(\xi_{n} > 0\) . Show that C is a pointed convex cone such that \(C \cap (-C) = \{0\}\) and \(C \cup (-C) = E\) ; deduce that C is the set of elements of E that are \(\geqslant 0\) for an order structure that is compatible with the vector space structure of E and for which E is linearly ordered. Show that on this ordered vector space the only linear positive form is 0.
\item
  Let E be an affine space of dimension \(\geqslant 2\) and let f be a bijection of E on itself; show that if the image under f of every convex subset of E is also a convex set, then f is an affine linear mapping (consider the inverse mapping of f and note that a closed segment is the intersection of the convex sets which contain its extremities cf.~A, II, § 9, exerc. 7).
\item
  Give an example of two convex sets \(\mathbf{A} \subset \mathbf{R}\) , \(\mathbf{B} \subset \mathbf{R}^2\) , such that the image of the convex set \(\mathbf{A} \times \mathbf{B}\) under the bilinear mapping \((\lambda, x) \mapsto \lambda x\) of \(\mathbf{R} \times \mathbf{R}^2\) in \(\mathbf{R}^2\) is not convex.
\item
  Let \((\mathbf{A}_i)_{1 \leqslant i \leqslant p}\) be a finite family of convex subsets of a vector space \(E\) ; let \(\mathbf{W}_i\) be the subspace obtained by a translation from the affine linear variety generated by \(\mathbf{A}_i\) ( \(1 \leqslant i \leqslant p\) ). If \(\mathbf{W} = \sum_{i=1}^{p} \mathbf{W}_i\) , show that the affine linear variety generated by the convex set \(\sum_{i=1}^{p} \lambda_i \mathbf{A}_i\) (where \(\lambda_i\) are non-zero numbers) is obtained from \(\mathbf{W}\) by a translation.
\end{enumerate}

¶ 9) Let A be a subset of the space \(\mathbf{R}^n\) .

\begin{enumerate}
\def\labelenumi{\alph{enumi})}
\tightlist
\item
  Show that the convex envelope of A is identical with the set of points \(\sum_{i=0}^{n} \lambda_i x_i\) , where \(x_i \in \mathbf{A}\) , \(\lambda_i \geqslant 0\) for \(0 \leqslant i < n\) , and \(\sum_{i=0}^{n} \lambda_i = 1\) . (Establish the following lemma; if \(p + 1\) points \(x_i\) ( \(0 \leqslant i \leqslant p\) ) form an affinely dependent system (that is to say there exists a relation \(\sum_{i=0}^{p} \beta_i x_i = 0\) where the \(\beta_i\) are not all zero and \(\sum_{i=0}^{p} \beta_i = 0\) ) and if \(x = \sum_{i=0}^{p} \alpha_i x_i\) , where the \(\alpha_i\) are \(\geqslant 0\) and \(\sum_{i=0}^{p} \alpha_i = 1\) , then there exists an index \(k \leqslant p\) and \(p\) numbers \(\gamma_i\) ( \(0 \leqslant i \leqslant p\) , \(i \neq k\) ) such that \(\gamma_i \geqslant 0\) for all \(i\) , \(\sum_{i \neq k} \gamma_i = 1\) and \(x = \sum_{i \neq k} \gamma_i x_i\) ; for this compare those of the \(\alpha_i / \beta_i\) that are defined.) b) Let \(a\) be a point of the convex envelope of A which does not belong to the convex envelope of any subset of A with at most \(n\) points. Show then that A contains at least \(n + 1\) connected components. (We can suppose that \(a = 0\) ; let \((b_i)_{0 \leqslant i \leqslant n}\) be a family of \(n + 1\) affinely independent points of A such that 0 belongs to the convex envelope of the \(b_i\) (cf.~a)). For each index \(i\) , let \(C_i\) be the pointed convex cone of vertex 0 generated by the \(b_j\) with indices \(j \neq i\) ; show that A does not meet the frontier of any of the cones - \(C_i\) .) c) If C is a pointed cone with vertex 0 in \(\mathbf{R}^n\) , show that the convex envelope of C is the set of points \(\sum_{i=1}^{n} x_i\) , where \(x_i \in C\) for \(1 \leqslant i \leqslant n\) .
\end{enumerate}

¶ 10) Let C be the convex envelope of a subset A of \(\mathbf{R}^n\) , and let \(a\) be an interior point of C. Show that there exist \(2n\) points \(x_i \in \mathrm{A}\) ( \(1 \leqslant i \leqslant 2n\) ) such that \(a\) is interior to \(\mathrm{C}_0\) , the convex envelope of the \(x_i\) . (Suppose \(a = 0\) , and argue by induction on \(n\) , noting that by exerc. \(9a\) ) there exists a set of \(k + 1\) points \(y_j\) of A ( \(0 \leqslant j \leqslant k\) , \(1 \leqslant k \leqslant n\) ), affinely independent and such that, if V is the affine linear variety generated by the \(y_j\) , then \(0 \in \mathrm{V}\) and, relative to V, 0 is interior to the convex envelope of the set of the \(y_j\) . Then project C on E/V and show that 0 is interior to this projection relative to E/V). Show that in the above statement \(2n\) cannot be replaced by \(2n - 1\) .

\begin{enumerate}
\def\labelenumi{\arabic{enumi})}
\setcounter{enumi}{10}
\tightlist
\item
  \begin{enumerate}
  \def\labelenumii{\alph{enumii})}
  \tightlist
  \item
    Show that, in the space \(\mathbf{R}^n\) , every convex set \(\mathbf{A}\) of dimension \(n\) contains at least one interior point (consider an affinely independent system of \(n + 1\) points of \(\mathbf{A}\) ). Deduce that if \(\mathbf{A}\) is everywhere dense in \(\mathbf{R}^n\) then \(\mathbf{A} = \mathbf{R}^n\) .
  \end{enumerate}
\end{enumerate}

\begin{enumerate}
\def\labelenumi{\alph{enumi})}
\setcounter{enumi}{1}
\item
  Let E be the normed space \(l^1(\mathbf{N})\) of absolutely convergent series of real numbers \(x = (\xi_n)\) (I, p.~4); show that the set P of \(x\) , such that \(\xi_n \geqslant 0\) for every index \(n\) , is a proper convex cone, which generates E but does not contain any interior point.
\item
  Let E be a Hausdorff topological vector space on which there exists a non-continuous linear form \(f\) (cf.~II, p.~86, exerc. 17, a)). Show that the sets A and B defined by the relations \(f(x) \geqslant 0\) , \(f(x) < 0\) are convex, non-empty, complementary, everywhere dense and that each of them generates E (algebraically).
\end{enumerate}

\begin{enumerate}
\def\labelenumi{\arabic{enumi})}
\setcounter{enumi}{11}
\item
  Show that in the space \(R^{n}\) , a necessary and sufficient condition that a convex set should be closed, is that its intersection with every straight line should be closed (cf.~II, p.~74, exerc. 5).
\item
  Show that in the space \(\mathbf{R}^n\) , every non-empty open convex set is homeomorphic to \(\mathbf{R}^n\) (use exerc. 12 of GT, VI, § 2).
\item
  Let A be a non-empty closed convex subset of E, a Hausdorff topological vector space.
\end{enumerate}

\begin{enumerate}
\def\labelenumi{\alph{enumi})}
\item
  Show that, for every \(a \in \mathbf{A}\) , the set \(\bigcap_{\lambda > 0} \lambda(\mathbf{A} - a)\) is a closed convex cone in \(\mathbf{E}\) , with vertex 0, independent of \(a\) . It is called the asymptotic cone of \(A\) and written \(C_A\) . For every \(a \in A\) , the set \(a + C_A\) is the union of \(\{a\}\) and those open half lines that are contained in \(A\) and have \(a\) as an end point.
\item
  If \(x, y\) are two points of \(\mathbf{A}\) such that \((x + \mathrm{C}_{\mathbf{A}}) \cap (y + \mathrm{C}_{\mathbf{A}})\) is a cone whose vertex \(z \in \mathbf{A}\) , then this cone is necessarily \(z + \mathrm{C}_{\mathbf{A}}\) .
\item
  If B is a second closed convex subset of E such that \(A \cap B \neq \varnothing\) , then \(C_{A \cap B} = C_{A} \cap C_{B}\) . d) Let \(V_{A}\) be the largest vector subspace (necessarily closed in E) which is contained in \(C_{A}\) . Show that if \(\phi\) is the canonical homomorphism of E on \(E/V_{A}\) , then \(A = \phi^{-1}(A_{0})\) , where \(A_{0}\) is a closed convex set in \(E/V_{A}\) which does not contain any straight line.
\item
  In the Banach space \(\mathcal{B}(\mathbf{N})\) of bounded mappings of N in R (I, p.~4) give an example of a closed convex non-bounded set A, for which \(C_{A} = \{0\}\) and which is such that for every \(b \neq 0\) in E, there exists an integer k for which \((\mathbf{A} + kb) \cap \mathbf{A} = \varnothing\) .
\end{enumerate}

\begin{enumerate}
\def\labelenumi{\arabic{enumi})}
\setcounter{enumi}{14}
\tightlist
\item
  \begin{enumerate}
  \def\labelenumii{\alph{enumii})}
  \tightlist
  \item
    Let A be a closed convex subset of a Hausdorff topological vector space E. If for some point \(x_0 \in \mathbf{A}\) there exists a neighbourhood \(\mathbf{V}\) of \(x_0\) in \(\mathbf{E}\) such that \(\mathbf{V} \cap \mathbf{A}\) is compact, then show that \(\mathbf{A}\) is locally compact. Deduce that the closure in \(\mathbf{E}\) of a locally compact convex set is locally compact.
  \end{enumerate}
\end{enumerate}

\begin{enumerate}
\def\labelenumi{\alph{enumi})}
\setcounter{enumi}{1}
\tightlist
\item
  Let A be a closed convex set that is locally compact but not compact in E; show that the asymptotic cone (exerc. 14) is not the single point \(\{0\}\) .
\end{enumerate}

¶ 16) Let A, B be two closed convex subsets of a Hausdorff topological vector space E. Suppose further that B is locally compact and that \(C_A \cap C_B = \{0\}\) . Show that A - B is closed in E. (Let \(b \in B\) , and W be a closed neighbourhood of 0 in E such that \(B \cap (b + W)\) is compact. Let \(c \in \overline{A - B}\) ; for every neighbourhood V of 0 in E consider the set \(M_V\) of those \(y \in B\) such that \(A \cap (c + y + V) \neq \emptyset\) . Consider two cases according as to whether there exists a V for which \(M_V\) is relatively compact, or there does not exist such a V; in the second case, consider the filter base formed from the sets \(P_{V,n} = M_V \cap \complement (b + nW)\) where V varies in the set of closed neighbourhood of 0 in E and n varies in N; form the cone with vertex b generated by \(P_{V,n}\) and its intersection with the frontier of \(b + W\) ).

¶ 17) In a Hausdorff topological vector space E, a closed convex set A is said to be parabolic if, for every \(z \notin A\) , each half-line originating at \(z\) and contained in \(z + C_A\) meets A.

\begin{enumerate}
\def\labelenumi{\alph{enumi})}
\tightlist
\item
  Give an example of a parabolic convex set A in \(\mathbf{R}^2\) such that \(C_A\) is not just a single half-line. b) Let A be a closed convex set in E such that \(C_A \neq \{0\}\) , but such that A is not parabolic. Show that if \(z \notin A\) is such that \(z + C_A\) contains a half-line D with end point \(z\) which does not meet A, then neither the convex envelope of \(A \cup \{z\}\) nor the pointed cone with vertex \(z\) generated by A is closed in E.
\end{enumerate}

Further if \(\mathbf{D}'\) is the closed half-line originating at \(z\) and opposite to \(\mathbf{D}\) (so that \(\mathbf{D} = 2z - \mathbf{D}'\) ) then \(\mathbf{D}' + \mathbf{A}\) is not closed in \(\mathbf{E}\) , and there exists a plane \(\mathbf{P}\) containing \(\mathbf{D}\) and a closed convex set \(\mathbf{B} \subset \mathbf{P}\) , such that \(\mathbf{B} \cap \mathbf{A} = \varnothing\) but that the distance of \(\mathbf{B}\) from \(\mathbf{P} \cap \mathbf{A}\) (in any norm on \(\mathbf{P}\) ) is zero.

\begin{enumerate}
\def\labelenumi{\alph{enumi})}
\setcounter{enumi}{2}
\item
  In E, let A be a closed convex set that is locally compact and parabolic; show that if \(\mathbf{B} \subset \mathbf{E}\) , is closed and convex then \(\mathbf{A} - \mathbf{B}\) is closed (same method as exerc. 16).
\item
  Let A, \(A'\) be two closed convex subsets of E that are locally compact and parabolic; show that the convex envelope of \(A \cup A'\) is closed in E (same method as in exerc. 16). Give an example in \(\mathbf{R}^2\) where \(\mathbf{A}\) is parabolic, \(\mathbf{A}'\) is non-parabolic and the convex envelope of \(\mathbf{A} \cup \mathbf{A}'\) is not closed in \(\mathbf{R}^2\) .
\item
  In E, let A be a closed convex set that is locally compact and parabolic; show that for all \(z \notin \mathbf{A}\) , the pointed convex cone with vertex \(z\) , generated by A, is closed in E (use \(d\) ).
\end{enumerate}

\(f\) ) Let \(\mathbf{E}_1, \mathbf{E}_2\) be two Hausdorff topological spaces, \(\mathbf{A}_1\) (resp. \(\mathbf{A}_2\) ) a closed convex set in \(\mathbf{E}_1\) (resp. \(\mathbf{E}_2\) ) that is also parabolic. Show that the set \(\mathbf{A}_1 \times \mathbf{A}_2\) is parabolic in \(\mathbf{E}_1 \times \mathbf{E}_2\) .

* g) Show that a barrelled space of infinite dimension does not contain a parabolic closed convex set that is both locally compact and not compact.

\(h\) ) Let \(\mathrm{E}_0 = l^2 (\mathbf{N})\) , and in \(\dot{\mathrm{E}}_0\) , let \(\mathbf{K}\) be the set of points \(x = (\xi_n)\) such that \(|\xi_n| \leqslant 1 / (n + 1)\) for all \(n\) ; \(\mathbf{K}\) is compact. The cone E of vertex 0 generated by \(\mathbf{K}\) is a vector subspace of \(\mathrm{E}_0\) and \(\mathbf{K}\) is absorbent in E. Let \(p\) be the gauge of \(\mathbf{K}\) in E; it is a lower semi-continuous function. In the normed product space \(\mathrm{E} \times \mathbf{R}\) , show that the set A of points \((x, \zeta)\) such that \(\zeta \geqslant (p(x))^2\) is closed, convex, parabolic and locally compact. Show that \(\mathrm{A} + (-\mathrm{A})\) and the convex envelope of \(\mathrm{A} \cup (-\mathrm{A})\) are not locally compact. *

\begin{enumerate}
\def\labelenumi{\arabic{enumi})}
\setcounter{enumi}{17}
\item
  Let \(\mathbf{C}\) be a proper closed convex cone of vertex 0 in \(\mathbf{R}^n\) . Show that the complement of the set \(\mathbf{C} \cap \mathbf{S}_{n-1}\) on the sphere \(\mathbf{S}_{n-1}\) is homeomorphic to \(\mathbf{R}^{n-1}\) (make a stereographic projection from a point of \(\mathbf{C} \cap \mathbf{S}_{n-1}\) , and use exerc. 12 of GT, VI, § 2). If \(\mathbf{C}\) contains an interior point, show that \(\mathbf{C} \cap \mathbf{S}_{n-1}\) , is homeomorphic to the closed ball \(\mathbf{B}_{n-1}\) (same method).
\item
  \begin{enumerate}
  \def\labelenumii{\alph{enumii})}
  \tightlist
  \item
    Let \(\mathbf{A}\) be an unbounded closed convex set in \(\mathbf{R}^n\) , that does not contain any line, but does contain an interior point. Show that the frontier of \(\mathbf{A}\) is homeomorphic to \(\mathbf{R}^{n-1}\) (use exercs. 15, b) and 18).
  \end{enumerate}
\end{enumerate}

\begin{enumerate}
\def\labelenumi{\alph{enumi})}
\setcounter{enumi}{1}
\tightlist
\item
  In a Hausdorff topological vector space E, let A be a closed convex set that does not contain any line and is of dimension \(\geqslant 2\) . Show that the frontier of A is connected (use a) and GT, VI, § 2, exerc. 12).
\end{enumerate}

\begin{enumerate}
\def\labelenumi{\arabic{enumi})}
\setcounter{enumi}{19}
\tightlist
\item
  \begin{enumerate}
  \def\labelenumii{\alph{enumii})}
  \tightlist
  \item
    In a vector space E, let A be a convex set that generates E and meets every straight line in a set that is closed relative to the straight line. Show that the following conditions are equivalent:
  \end{enumerate}
\end{enumerate}

\(\alpha)\) There exists a line D such that D meets A in a compact segment that is not empty.

β) There exists a line D such that every line parallel to D meets A in a compact segment.

\(\gamma)\) A is distinct from E and is not a half-space determined by a hyperplane of \(\tilde{\mathbf{E}}\) .

(To show that \((\gamma) \Rightarrow (\alpha)\) use exerc. 14, \(d\) ) of II, p.~67, and reduce to the case \(\mathbf{E} = \mathbf{R}^2\) .)

\begin{enumerate}
\def\labelenumi{\alph{enumi})}
\setcounter{enumi}{1}
\tightlist
\item
  In a Hausdorff topological vector space E, let A be a closed convex set which contains an interior point. Show that if the frontier of A is a non-empty linear variety, then A is a closed half-space (use exerc. 14, d) of II, p.~67, to show that the frontier of A is necessarily a hyperplane, then apply a)).
\end{enumerate}

¶ 21) a) Let \(\mathrm{A}_i\) ( \(1 \leqslant i \leqslant r\) ) \(r > n + 1\) , be a family of convex subsets of \(\mathbf{R}^n\) such that any \(r - 1\) of the \(\mathrm{A}_i\) have a non-empty intersection; show that the \(r\) sets \(\mathrm{A}_i\) have a non-empty intersection (Helly's theorem). (Let \(x_i\) be a point of the intersection of the \(\mathrm{A}_j\) with indexes \(j \neq i\) ;

there exist \(r\) numbers \(\lambda_{i}\) which are not all zero and are such that \(\sum_{i=1}^{r} \lambda_{i} = 0\) and \(\sum_{i=1}^{r} \lambda_{i} x_{i} = 0\) ;

in this last equation take to one side those terms with \(\lambda_{i} \geqslant 0\) and to the other those with \(\lambda_{i} < 0\) .)

\begin{enumerate}
\def\labelenumi{\alph{enumi})}
\setcounter{enumi}{1}
\item
  Given a family of compact convex sets in \(\mathbf{R}^n\) , show that the intersection of all the sets of the family is non-empty if the intersection of any selection of \(n + 1\) sets of the family is non-empty.
\item
  In \(\mathbf{R}^n\) , let \(\mathbf{K}\) be a convex set and \((\mathbf{A}_i)_{1 \leqslant i \leqslant r}\) be a family of \(r > n + 1\) convex sets. Suppose that for every selection of \(n + 1\) indices \((i_k)\) each less than or equal to \(r\) there exists \(a \in \mathbf{R}^n\) such that \(a + \mathbf{K}\) contains each of the \(A_{i_k}\) . Show that then there exists \(b \in \mathbf{R}^n\) such that \(b + \mathbf{K}\) contains all the \(\mathbf{A}_i\) . Show that similar results hold if « contains » is replaced by « is contained in » or by « meets in a non-empty set ». (For each index \(i\) , consider the set \(\mathbf{C}_i\) of the \(x \in \mathbf{R}^n\) for which \(x + \mathbf{K} \supset \mathbf{A}_i\) (or \(x + \mathbf{K} \subset \mathbf{A}_i\) , or \((x + \mathbf{K}) \cap \mathbf{A}_i \neq \emptyset\) )). Generalize to any family of compact convex sets of \(\mathbf{R}^n\) .
\end{enumerate}

\begin{enumerate}
\def\labelenumi{\arabic{enumi})}
\setcounter{enumi}{21}
\item
  In \(\mathbf{R}^2\) consider a set of \(2m\) points of the form \((a_i, b_i')\) , \((a_i, b_i'')\) where \(b_i' \leqslant b_i''\) for \(1 \leqslant i \leqslant m\) . Let \(n\) be an integer \(< m - 2\) . In order that there should exist a polynomial \(\mathrm{P}(x)\) of degree \(\leqslant n\) such that \(b_i' \leqslant \mathrm{P}(a_i) \leqslant b_i''\) for \(1 \leqslant i \leqslant m\) , it is sufficient that, for every family \((i_k)_{1 \leqslant k \leqslant n + 2}\) of \(n + 2\) indices \(i\) , there exists a polynomial \(\mathrm{Q}(x)\) of degree \(\leqslant n\) such that \(b_{i_k}' \leqslant \mathrm{Q}(a_{i_k}) \leqslant b_{i_k}''\) for every integer \(k\) such that \(1 \leqslant k \leqslant n + 2\) . (Use exerc. 21, a).)
\item
  Show that in a topological vector space, the convex envelope of an open set is an open set.
\item
  Let \(\mathbf{M}\) be an everywhere dense convex set in a topological vector space \(\mathrm{E}\) (cf.~II, p.~66, exerc. 11, c)); show that, for every closed hyperplane \(\mathbf{H}\) in \(\mathbf{E}\) , the set \(\mathbf{H} \cap \mathbf{M}\) is dense in \(\mathbf{H}\) (for every point \(x_0 \in H\) , and every balanced neighbourhood \(\mathbf{V}\) of 0 in \(\mathbf{E}\) , consider the intersections of \(x_0 + \mathbf{V}\) and the two open half-spaces determined by \(\mathbf{H}\) , and deduce that \(x_0 + \mathbf{V} + \mathbf{V}\) meets \(\mathbf{H} \cap \mathbf{M}\) ).
\item
  \begin{enumerate}
  \def\labelenumii{\alph{enumii})}
  \tightlist
  \item
    Show that, in a topological vector space, every convex set with an interior point, is such that its frontier is nowhere dense (use prop. 16 of II, p.~14).
  \end{enumerate}
\end{enumerate}

\begin{enumerate}
\def\labelenumi{\alph{enumi})}
\setcounter{enumi}{1}
\tightlist
\item
  In a Hausdorff topological vector space E, let A be a closed convex set with an interior point, and let H be a closed hyperplane that contains an interior point of A. Show that the intersection of H and of the frontier F of A is a set which is nowhere dense relative to F (to show that in every neighbourhood of a point of \(\mathbf{H} \cap \mathbf{F}\) there exist points of F not in H, reduce to the case when E is of dimension 2).
\end{enumerate}

¶ 26) In a Hausdorff topological vector space E, let A be a connected closed set with the following property : for every \(x \in A\) , there exists a closed neighbourhood V of \(x\) in E such that V ∩ A is convex. Show that A is convex. For this establish the following statements.

\begin{enumerate}
\def\labelenumi{\alph{enumi})}
\item
  Show that any two points of A can be joined by a broken line in A (same method as GT, VI, § 1, exerc. 6).
\item
  Show that, if two points in A can be joined by a broken line in A with n > 1 segments, then they can also be joined in A by a broken line with n - 1 segments. (Induction on n reduces to the case n = 2 which is equivalent to taking \(R^{2}\) as E; then let T be a triangle with vertices a, b, c such that the closed segments ac, bc are contained in A, but the closed segment ab is not; consider a point of the closure of the intersection of \(\complement A\) and the interior of T that is farthest from the line ab, and show that the existence of such a point contradicts the hypothesis.)
\end{enumerate}

¶ 27) a) Let B be a non-empty closed convex subset of E, a Hausdorff topological vector space, and let X be a non-empty compact set in E. Show that if A is a subset of E such that \(A + X \subset B + X\) , then \(A \subset B\) (if \(a \in A\) , consider a sequence \((x_{n})\) of points of X defined inductively by the relation \(a + x_{n} = b_{n} + x_{n+1}\) , where \(b_{n} \in B\) ). Deduce that, if A, B are two non-empty subsets of E, using the distance in E and the procedure of GT, IX, § 2, exerc. 6. Show \(A + X = B + X\) implies the relation A = B.

\begin{enumerate}
\def\labelenumi{\alph{enumi})}
\setcounter{enumi}{1}
\item
  Let E be a normed space, \(\sigma\) the distance function defined on the set \(\mathfrak{F}(E)\) of closed non-empty subsets of E, using the distance in E and the procedure of TG, IX, p.~91, exerc. 6. Show that if A, B, C are three non-empty compact convex sets in E then \(\sigma(A+C, B+C)=\sigma(A, B)\) (if \(S_{\lambda}\) is the ball defined by \(\|x\|\leqslant\lambda\) , note that \(A+S_{\lambda}\) and \(B+S_{\lambda}\) are closed convex sets and use a)).
\item
  Deduce from a) and b) that the set \(\Re(\mathrm{E})\) of non-empty, compact, convex subsets of a normed space E, with the distance \(\sigma\) , can be identified with a cone in a normed space of which the laws of composition induce on \(\Re(\mathrm{E})\) the laws \((\mathrm{A}, \mathrm{B}) \to \mathrm{A} + \mathrm{B}\) and \((\lambda, \mathrm{A}) \to \lambda\mathrm{A}\) .
\end{enumerate}

\begin{enumerate}
\def\labelenumi{\arabic{enumi})}
\setcounter{enumi}{27}
\tightlist
\item
  Let \(f\) be a convex function defined over the convex subset \(A\) of a vector space \(E\) .
\end{enumerate}

\begin{enumerate}
\def\labelenumi{\alph{enumi})}
\item
  Show that if \(\mathbf{A}\) is absorbent and \(f\) is non-constant then \(f\) cannot attain its upper bound in \(\mathbf{A}\) at the point 0.
\item
  Show that the subset of points of A, at which \(f\) attains its lower bound in A, is convex.
\end{enumerate}

¶ 29) Let E be a Hausdorff topological vector space, and C be a non-empty open convex non-pointed cone with vertex 0, in E. A convex neighbourhood of 0 in E is denoted by V. If f is a convex function that is defined and bounded above in C ∩ V, show that \(f(x)\) tends to a finite limit as x tends to 0, where \(x \in C \cap V\) . (Let \(\beta = \limsup_{x \to 0, x \in C \cap V} f(x)\) ; obtain a contradiction,

supposing that for some \(\alpha > 0\) , and every neighbourhood \(\mathbf{W}\) of \(0\) , there exists a point \(y \in \mathbf{C} \cap \mathbf{V} \cap \mathbf{W}\) such that \(f(y) < \beta - \alpha\) . Show that there exists \(a \in \mathbf{C} \cap \mathbf{V}\) such that \(f(\rho a) \geqslant \beta - \frac{1}{2}\alpha\) for \(0 < \rho \leqslant 1\) ; deduce that, on a line joining a point of the form \(\rho a\) ( \(\rho\) sufficiently small) to a point \(y\) in \(\mathbf{C} \cap \mathbf{V}\) such that \(f(y) < \beta - \alpha\) , and that is sufficiently close to \(0\) , there exist points of \(\mathbf{C} \cap \mathbf{V}\) where \(f\) is arbitrarily large.)

\begin{enumerate}
\def\labelenumi{\arabic{enumi})}
\setcounter{enumi}{29}
\tightlist
\item
  \begin{enumerate}
  \def\labelenumii{\alph{enumii})}
  \tightlist
  \item
    Give an example of a convex function that is defined over a compact convex subset K of \(R^{2}\) , that is bounded and lower semi-continuous in K, but is not continuous at a point of the frontier of K (consider the gauge of a disc of which 0 is a frontier point).
  \end{enumerate}
\end{enumerate}

\begin{enumerate}
\def\labelenumi{\alph{enumi})}
\setcounter{enumi}{1}
\item
  Deduce from \(a\) an example of a convex function defined in an open half-plane \(\mathbf{D}\) of \(\mathbb{R}^2\) , not bounded above in \(\mathbf{D}\) and not tending to a limit at a frontier point of \(\mathbf{D}\) .
\item
  Deduce from a) an example of a convex lower semi-continuous function defined over a compact convex subset A of \(R^{2}\) but not bounded above in A. (Take for A the set of points \((\xi, \eta)\) such that \(\xi^{4} \leqslant \eta \leqslant 1\) in \(R^{2}\) .)
\end{enumerate}

¶ 31) Let \(x_0\) be a point in the closure of A, a non-empty convex subset of a Hausdorff topological vector space E. Let \(f\) be a convex function defined over A. Use \(\mathfrak{D}\) to denote the set of closed half lines D which originate at \(x_0\) , for which \(A \cap D\) contains an open segment with end point \(x_0\) . The union C of the half lines \(D \in \mathfrak{D}\) is a convex cone with vertex \(x_0\) .

\begin{enumerate}
\def\labelenumi{\alph{enumi})}
\item
  Show that, for each fixed \(\mathbf{D} \in \mathfrak{D}\) , as \(x\) tends to \(x_0\) such that \(x \in \mathbf{D} \cap \mathbf{A}\) and \(x \neq x_0\) , either \(f(x)\) tends to a finite limit or to \(+\infty\) .
\item
  Let \(\mathfrak{I}\) be the subset of those \(D \in \mathfrak{D}\) , for which the limit of \(f(x)\) in \(a\) is \(+\infty\) ; if \(x_0 \in A\) then \(\mathfrak{I}\) is empty. Show that \(\mathfrak{I}\) cannot contain two opposite half lines; if \(D\) and \(D'\) are two distinct half lines in \(\mathfrak{I}\) and \(P\) is the plane determined by \(D\) and \(D'\) then, either, every half line \(D''\) of \(\mathfrak{D}\) in \(P\) belongs to \(\mathfrak{I}\) , or \(D\) and \(D'\) are the only two half lines of \(\mathfrak{I}\) lying in \(P\) . Deduce that if \(\mathfrak{I} \neq \mathfrak{D}\) , then no half line \(D \in \mathfrak{I}\) contains an internal point (II, p.~26) of the cone \(C\) relative to the vector subspace generated by \(C\) .
\item
  Let \(\mathfrak{F}\) be the set of half lines in \(\mathfrak{D}\) that are not in \(\mathfrak{I}\) . Show that the union of the half lines of \(\mathfrak{F}\) is a convex cone, and for each half line \(\mathrm{D} \in \mathfrak{F}\) the limit of \(f(x)\) defined in \(a\) ) is independent of \(\mathrm{D}\) (use exerc. 29 above); further if \(x_0 \in \mathrm{A}\) this limit is \(\leqslant f(x_0)\) , and it is equal to \(f(x_0)\) when \(\mathfrak{F}\) contains two opposite half lines.
\end{enumerate}

\(d\) ) Let \(f\) be a non-continuous linear form over \(\mathbf{E}\) (cf.~II, p.~86, exerc. 17, \(a\) ) and take \(\mathbf{A} = \mathbf{E}\) ; show that every closed half line, originating at \(x_0\) , belongs to \(\mathfrak{F}\) , but that

\[
\lim _ {x \to x _ {0}} \inf f (x) = - \infty \text { and } \lim _ {x \to x _ {0}} \sup f (x) = + \infty
\]

(use prop. 21 of II, p.~18).

\begin{enumerate}
\def\labelenumi{\arabic{enumi})}
\setcounter{enumi}{31}
\item
  Let \(\mathbf{K}\) be a compact convex set in a Hausdorff topological vector space \(\mathbf{E}\) and let \(f\) be an upper semi-continuous convex function defined over \(\mathbf{K}\) . Show that \(f\) is bounded over \(\mathbf{K}\) . (Observe first that \(f\) is bounded above in \(\mathbf{K}\) ; if \(f\) is not bounded below show that \(\liminf_{y\to x,y\neq x}f(y) = -\infty\) for every point \(x\in \mathbf{K}\) , and that this contradicts Baire's theorem.) Give an example where \(f\) is not continuous.
\item
  Let E be a finite dimensional Hausdorff topological vector space, and let K be a compact convex subset of E. Show that every convex function defined over K is bounded below in K (compare exerc. 31, d)).
\item
  Let U, V be two open convex sets in a Hausdorff topological vector space E such that \(\overline{V} \subset U\) and that U does not contain any half line. Let F be a set of convex functions defined in \(\overline{U}\) , uniformly bounded above on the frontier of U and uniformly bounded below on the frontier of V. Show that F is equicontinuous.
\item
  Let U be a non-empty open convex set in \(R^{n}\) and F be a set of convex functions defined over U. Let \(\Phi\) be a filter on F that converges pointwise in U to a finite function \(f_{0}\) ; show that \(\Phi\) converges uniformly to \(f_{0}\) in every compact subset of U (use exerc. 34).
\item
  Let A be a compact convex set in \(\mathbf{R}^n\) and B its projection on the subspace \(\mathbf{R}^{n-1}\) (identified as the hyperplane with equation \(\xi_n = 0\) ). Show that there exist two convex functions \(f_1, f_2\) defined over B, such that A is identical with the set of points \((x, \zeta)\) of \(\mathbf{R}^n\) where \(x \in \mathbf{B}, y \in \mathbf{R}\) and \(f_1(x) \leqslant \zeta \leqslant -f_2(x)\) .
\item
  Let E be a vector space; in order that a convex set F of \(E \times R\) should be formed of pairs \((x, \zeta)\) such that \(f(x) \leqslant \zeta\) (resp. \(f(x) < \zeta\) ) for a convex function f defined over a convex subset X of E, it is necessary and sufficient that the projection of F on E should be identical with X and that, for all \(x \in X\) , the set \(F(x)\) of F that projects onto x should be a closed (resp. open) interval unlimited to the right (i.e.~not bounded above).
\item
  Let \(\mathbf{X}\) be a convex set of an affine space \(\mathbf{E}\) and \(p\) an affine linear mapping of \(\mathbf{E}\) in a second affine linear space \(\mathbf{E}_1\) . Write \(\mathbf{X}_1 = p(\mathbf{X})\) . For every real-valued function \(f\) defined in \(\mathbf{X}\) and every \(x_1 \in \mathbf{X}_1\) , let
\end{enumerate}

\[
f _ {1} (x _ {1}) = \inf _ {p (x) = x _ {1}} f (x).
\]

Show that if \(f\) is convex and if \(f_{1}(x_{1}) > -\infty\) for all \(x_{1} \in X_{1}\) , then \(f_{1}\) is a convex function.

\begin{enumerate}
\def\labelenumi{\arabic{enumi})}
\setcounter{enumi}{38}
\tightlist
\item
  Let \(\mathbf{E}\) be a finite dimensional Hausdorff topological vector space.
\end{enumerate}

\begin{enumerate}
\def\labelenumi{\alph{enumi})}
\item
  Let \(\mathfrak{F}(\mathrm{E})\) be the family of closed non-empty sets of E, carrying the uniform structure deduced from the uniform structure of E by the procedure of GT, II, § 1, exerc. 5, a). Show that, the set \(\mathfrak{C}(\mathrm{E})\) of non-empty closed convex sets of E, is closed in the space \(\mathfrak{F}(\mathrm{E})\) . Deduce that if K is a compact set in E, the set of non-empty closed convex sets in E that are contained in K, is a compact set in \(\mathfrak{C}(\mathrm{E})\) (cf.~GT, § 4, exerc. 11).
\item
  Let \(\mathfrak{R}_0(\mathrm{E})\) be the set of compact convex subsets of \(\mathrm{E}\) that contain 0 as an interior point. For every set \(\mathrm{A} \in \mathfrak{R}_0(\mathrm{E})\) , let \(p_{\mathrm{A}}\) be the gauge of \(\mathrm{A}\) (II, p.~20). Show that \(\mathrm{A} \mapsto p_{\mathrm{A}}\) is an isomorphism of the uniform subspace \(\mathfrak{R}_0(\mathrm{E})\) of \(\mathfrak{C}(\mathrm{E})\) on a subspace of the space \(\mathcal{C}_{c}(\mathrm{E};\mathbf{R})\) of continuous real valued functions in \(\mathrm{E}\) , carrying the uniform structure of compact convergence (GT, X, § 1.6).
\end{enumerate}

\begin{enumerate}
\def\labelenumi{\arabic{enumi})}
\setcounter{enumi}{39}
\item
  In a topological vector space E, let U be a convex neighbourhood of \(x_{0}\) and let f be a real-valued continuous convex function in U. Show that there exists a convex neighbourhood \(V \subset U\) of \(x_{0}\) and a convex continuous function \(f_{1}\) in E such that \(f_{1}|V = f|V\) .
\item
  Let H be a hyperplane in a vector space E that does not contain 0 and let S be a convex set contained in H.
\end{enumerate}

\begin{enumerate}
\def\labelenumi{\alph{enumi})}
\item
  Suppose that the intersection of S with each line in H is a compact segment. Let \(a, b\) be two distinct points of E such that there exist two numbers \(\lambda > 0\) , \(\mu > 0\) for which \(b + \mu S \subset a + \lambda S\) ; show that if \(c\) is the point where the line joining \(a\) and \(b\) meets the hyperplane H' parallel to H which contains \(a + \lambda S\) , then \(c \in b + \mu S\) and \(b + \mu S\) is the image of \(a + \lambda S\) by a homothety of centre \(c\) transforming \(a\) into \(b\) . (Reduce to the case where E is of dimension 2.)
\item
  With the same hypotheses on S, let a, b be two distinct points of E and suppose that there exists a point \(c \in E\) and three numbers \(\lambda > 0\) , \(\mu > 0\) , \(v \geqslant 0\) such that \((a + \lambda S) \cap (b + \mu S) = c + vS\) . Show that if A (resp. B) is the cone with vertex a (resp. b) generated by \(a + \lambda S\) (resp. \(b + \mu S\) ), and \(H''\) the hyperplane parallel to H passing through c, then \(H'' \cap A \cap B = \{c\}\) (use a)). c) Suppose that H is the affine linear variety generated by S. Let C be the cone with vertex 0 generated by S. Show that the following two conditions are equivalent:
\end{enumerate}

α) E is a lattice for the order on E of which C is the set of elements ≥ 0.

β) For any points \(x, y\) of \(E\) and numbers \(\lambda > 0\) , \(\mu > 0\) such that the set \((x + \lambda S) \cap (y + \mu S)\) is not empty, there exists \(z \in E\) and \(\nu \geqslant 0\) such that this set is \(z + \nu S\) .

(To prove that \(\alpha\) ) implies \(\beta\) ), reduce to the case \(y = 0\) and use the fact that if \((s_i)\) is a finite family of points of S and \((\lambda_i)\) is a family of real numbers such that \(\sum_{i} \lambda_i s_i = 0\) , then \(\sum_{i} \lambda_i = 0\) .

To prove that \(\beta)\) implies \(\alpha)\) use \(b)\) .

When S satisfies the equivalent conditions \(\alpha)\) and \(\beta)\) , we say that S is a simplex in E. When E is of finite dimension, the convex envelope of a finite set of points affinely independent in H and generating H is a simplex * (the converse is also true : cf.~INT, II, 2nd ed., § 2, exerc. 7)).

\begin{enumerate}
\def\labelenumi{\arabic{enumi})}
\setcounter{enumi}{41}
\item
  Generalise the prop. 18 of II, p.~16 to the case of an ordered set E with a Hausdorff topology for which the intervals \(\{a, \to (\text{and}) \gets, a\}\) are closed for all \(a \in E\) .
\item
  Let K be a complete valued division ring of which the absolute value is an ultrametric. In a left vector space E on K, we say that a set A is ultraconvex if the relations \(x \in A\) , \(y \in A\) , \(|\lambda| \leqslant 1\) , \(|\mu| \leqslant 1\) imply \(\lambda x + \mu y \in A\) .
\end{enumerate}

\begin{enumerate}
\def\labelenumi{\alph{enumi})}
\tightlist
\item
  Generalize the prop. 1, 2, 5, 6, 7 of II, p.~8 and p.~9. Show that the smallest ultraconvex set containing a given set \(\mathbf{M}\) is the set of linear combinations \(\sum_{\iota} \lambda_{\iota} x_{\iota}\) , where \(x_{\iota} \in \mathbf{M}\) and \(|\lambda_{\iota}| \leqslant 1\)
\end{enumerate}

for all \(\iota\) .

\begin{enumerate}
\def\labelenumi{\alph{enumi})}
\setcounter{enumi}{1}
\item
  Suppose that E is a topological vector space over K. Show that the closure of an ultraconvex set is ultraconvex, and that an ultraconvex set with a non-empty interior is open.
\item
  Let A be an absorbent and ultraconvex set in E. Show that if, for all \(x \in E\) , we put \(p(x) = \inf_{x \in \mathbb{P}^{\mathbf{A}}} |\rho|\) , then p is an ultra-semi-norm on E (II, p.~2). Generalize prop. 23 of II, p.~20,
\end{enumerate}

to the case where the absolute value of K is obtained from a discrete valuation.

\exercisegroup

\begin{enumerate}
\def\labelenumi{\arabic{enumi})}
\item
  Let P be a proper pointed convex cone, with vertex 0, in E a vector space over R and p be a semi-norm on E and V the set of points \(x \in E\) for which \(p(x) < 1\) . Let M be a vector subspace of E and f be a linear form on M. There exists a linear form g on E, which extends f and is such that it is \(\geqslant 0\) in P and \(|g(x)| \leqslant p(x)\) for all \(x \in E\) , if and only if, for all \(x \in M \cap (V + P)\) , we have \(f(x) > -1\) . (To see that the condition is sufficient consider a point \(x_0 \in M\) such that, \(f(x_0) = 1\) , the cone Q of vertex 0 generated by \(x_0 + V\) , and apply the cor. of the prop. 1 (II, p.~21) to the space E carrying the relation of preorder for which \(P + Q\) is the cone of elements \(\geqslant 0\) .)
\item
  For a set S let \(F = \mathcal{B}(S)\) be the Banach space of the real-valued bounded functions in S (I, p.~4) and let M be a vector subspace of a normed space E. Show that, for every continuous linear mapping f of M in F, there exists a continuous linear mapping g of E in F, that is an extension of f and such that \(\|g\| = \|f\|\) .
\item
  Let E be a vector space over R and let p be a sublinear function on E (II, p.~20). Let A be a convex set such that \(\inf p(y) > -\infty\) .
\end{enumerate}

\begin{enumerate}
\def\labelenumi{\alph{enumi})}
\tightlist
\item
  Show that the function
\end{enumerate}

\[
q (x) = \inf _ {z \in \mathrm{A}, t \geqslant 0} \left(p (x + t z) - t. \inf _ {y \in \mathrm{A}} p (y)\right)
\]

is a sublinear function on E such that \(-p(-x) \leqslant q(x) \leqslant p(x)\) . b) Show that there exists a linear form h on E such that \(h(x) \leqslant p(x)\) in E and that \(\inf_{y \in A} p(y) = \inf_{y \in A} h(y)\) (take h such that \(h(x) \leqslant q(x)\) ).

\begin{enumerate}
\def\labelenumi{\arabic{enumi})}
\setcounter{enumi}{3}
\tightlist
\item
  Let A be a non-empty set of E, a vector space over R, and p a sub-linear function on E. Let B be the set of \(z \in E\) such that \(\inf_{x \in A} p(x - z) \leqslant 0\) ; we have \(A \subset B\) and \(\inf_{x \in A} p(x) \leqslant p(z)\) for all \(z \in B\) ; from which \(\inf_{x \in A} p(x) = \inf_{z \in B} p(z)\) .
\end{enumerate}


\begin{enumerate}
\def\labelenumi{\alph{enumi})}
\item
  Show that the set of the \(y \in \mathbf{E}\) such that \(\inf_{z \in \mathbf{B}} p(z - y) \leqslant 0\) is the set \(\mathbf{B}\) .
\item
  Deduce from a) that the intersection of B and any affine line D in E is closed in D (show that, whatever the points a, b of E, the function \(t \rightarrow p(a + tb)\) is continuous in R).
\item
  Suppose that for each pair of points x, y of A there exists \(z \in A\) such that \(p(z - \frac{1}{2}(x + y)) \leqslant 0\) . Show that, for each pair of points u, v of B we have \(\frac{1}{2}(u + v) \in \mathbf{B}\) (write
\end{enumerate}

\[
z - \frac {1}{2} (u + v) = \left(z - \frac {1}{2} (x + y)\right) + \frac {1}{2} (x - u) + \frac {1}{2} (y - v)
\]

for \(x, y, z\) in A). Deduce that B is then convex (use \(b\) ).

\begin{enumerate}
\def\labelenumi{\alph{enumi})}
\setcounter{enumi}{3}
\tightlist
\item
  Under the hypotheses of \(c\) ), show that there exists a linear form \(h\) on \(\mathbf{E}\) such that \(h(x) \leqslant p(x)\) and that we have \(\inf_{y \in \mathbf{A}} p(y) = \inf_{y \in \mathbf{A}} h(y)\) (use \(c\) ) and exerc. 3).
\end{enumerate}

\begin{enumerate}
\def\labelenumi{\arabic{enumi})}
\setcounter{enumi}{4}
\item
  Let A be a non-empty subset of a vector space E over R and let p be a sublinear function on E. Suppose that, for every pair of points x, y of A, there exists \(z \in A\) such that \(p(z - (x + y)) \leqslant 0\) and that \(p(x) \geqslant 0\) for all \(x \in A\) . Show that there exists a linear form h on E such that \(h(x) \leqslant p(x)\) in E and that \(h(x) \geqslant 0\) for \(x \in A\) . (Apply exerc. 4, c) to the union of the \(\frac{1}{n}A\) for n (integer) \(\geqslant 1\) .)
\item
  \begin{enumerate}
  \def\labelenumii{\alph{enumii})}
  \tightlist
  \item
    Let H be a hyperplane in E a vector space over R and let p be a sublinear function on E. Let f be a linear form on H, such that \(f(y) \leqslant p(y)\) in H. Let a be a point of \(\complement H\), and let h be the linear form on E which extends f and is such that \(h(a) = \inf_{y \in H} (f(y + p(a - y)))\) . Then \(h(x) \leqslant p(x)\) in E. Show that for every linear form \(g\) on E extending \(f\) and such that \(g(x) \leqslant p(x)\) in E, we also have \(g(a) \leqslant h(a)\) .
  \end{enumerate}
\end{enumerate}

\begin{enumerate}
\def\labelenumi{\alph{enumi})}
\setcounter{enumi}{1}
\tightlist
\item
  Let V be a vector subspace of E and f a linear form on V such that \(f(y) \leqslant p(y)\) in V. Let S be a non-empty set of E. Show that there exists a linear form h on E that extends f, such that \(h(x) \leqslant p(x)\) in E and that there is no other linear form g on E extending f such that \(g(x) \leqslant p(x)\) in E that is distinct from h and such that \(g(x) \geqslant h(x)\) in S. (Consider the set F of pairs \((V', f')\) where \(V'\) is a vector subspace containing V and \(f'\) a linear form on \(V'\) extending f and such that \(f'(z) \leqslant p(z)\) in \(V'\) and further such that there is no other linear form \(f''\) on \(V'\) with the same properties and such that \(f''(z) \geqslant f'(z)\) in \(S \cap V'\) . Order F and use a) and th. 2 of S, III, § 2.4.)
\end{enumerate}

\begin{enumerate}
\def\labelenumi{\arabic{enumi})}
\setcounter{enumi}{6}
\tightlist
\item
  Let \(\mathrm{T}\) be a commutative monoid (A, I, § 2.1) carrying a preorder relation \(x \leqslant y\) such that if \(x \leqslant y\) then \(x + z \leqslant y + z\) for all \(z \in \mathrm{T}\) . A mapping \(f\) of \(\mathrm{T}\) in \(\mathbf{R} \cup \{-\infty\}\) is called additive (resp. subadditive, resp. superadditive) if we have
\end{enumerate}

\[
f (x + y) = f (x) + f (y) \quad (\text { resp. } f (x + y) \leqslant f (x) + f (y), \quad \text { resp. } f (x + y) \geqslant f (x) + f (y))
\]

for any \(x, y\) in \(\mathbf{T}\) .

\begin{enumerate}
\def\labelenumi{\alph{enumi})}
\item
  If \(g\) is subadditive and increasing in \(\mathbf{T}\) , then the function \(h(x) = \inf_{n > 0} g(nx) / n\) is subadditive and increasing; we have \(h \leqslant g\) and \(h(0) = 0\) if \(g(0) \geqslant 0\) .
\item
  Under the same hypotheses suppose that there exist two elements \(x_{1}, x_{2}\) of \(T\) and two real numbers \(\xi_{1}, \xi_{2}\) such that \(\xi_{1} < g(x_{1}), \xi_{2} < g(x_{2})\) and \(g(x_{1} + x_{2}) < \xi_{1} + \xi_{2}\) . Let \(y_{1}, y_{2}, z_{1}, z_{2}\) be four elements of \(T\) , let \(n_{1}, n_{2}\) be two integers \(\geqslant 0\) and let \(\alpha_{1}, \alpha_{2}\) be two real numbers such that
\end{enumerate}

\[
\begin{array}{l} n _ {1} \xi_ {1} + g (z _ {1}) <   \alpha_ {1}, y _ {1} \leqslant n _ {1} x _ {1} + z _ {1} \\ n _ {2} \xi_ {2} + g (z _ {2}) <   \alpha_ {2}, y _ {2} \leqslant n _ {2} x _ {2} + z _ {2}. \end{array}
\]

Show that then \(g(n_2y_1 + n_1y_2) < n_2\alpha_1 + n_1\alpha_2\) .

\begin{enumerate}
\def\labelenumi{\alph{enumi})}
\setcounter{enumi}{2}
\tightlist
\item
  Let \(\omega\) be a superadditive function on T such that \(\omega(0) = 0\) , and let \(\Omega\) be an increasing subadditive function on T such that \(\omega(x) \leqslant \Omega(x)\) in T. Show that there exists an increasing additive function \(f\) on T such that \(\omega(x) \leqslant f(x) \leqslant \Omega(x)\) in T. (Remark that the set of increasing subadditive functions \(g\) on T such that \(\omega(x) \leqslant g(x) \leqslant \Omega(x)\) in T is non-empty and inductive for the relation \(\geqslant\) , and take a minimal element of this set for \(f\) ; show using \(a\) ) that \(f(0) = 0\) . To show that there cannot exist pairs of elements of T, \((x_1, x_2)\) such that \(f(x_1 + x_2) < f(x_1) + f(x_2)\) remark that if \(\xi_j \in \mathbf{R}\) , and \(h_j(x) = \inf(n\xi_j + f(y))\) where \(n\) varies in the set of integers \(\geqslant 0\) and \(y\) in the set of elements of T such that \(x \leqslant nx_j + y\) , then \(h_j\) is increasing and subadditive in T ( \(j = 1, 2\) ), \(h_j(x_j) \leqslant \xi_j\) and \(h_j(x) \leqslant f(x)\) for all \(x \in T\) . Then use the definition of \(f\) and part \(b\) ) to obtain a contradiction.)
\end{enumerate}

¶ * 8) a) Let K be a non-discrete valued division ring of which the absolute value is an ultrametric, non linearly compact(cf.~CA, III, § 2, exerc. 15); then there exists a well ordered set I of numbers > 0 and a family (B(ρ)) \(_{ρ \in I}\) of closed balls in K such that the relation ρ < ρ' implies B(ρ) ⊂ B(ρ'), that B(ρ) has radius ρ and that the intersection of the B(ρ) is empty (CA, VI, § 5, exerc. 5). For every x ∈ K, there exists ρ ∈ I such that x ∉ B(ρ); show that the number φ(x) = \(|x-y|\) for a \(y \in \mathbf{B}(\rho)\) depends neither on \(y \in \mathbf{B}(\rho)\) , nor on \(\rho \in \mathbf{I}\) such that \(x \notin \mathbf{B}(\rho)\) . If \(\rho \in \mathbf{I}\) is such that \(x \in \mathbf{B}(\rho)\) then \(\phi(x) \leqslant \rho\) . This being so, for \((x_1, x_2) \in \mathbf{K}^2\) , put \(\| (x_1, x_2) \| = |x_1|\) if \(x_2 = 0\) , and \(\| (x_1, x_2) \| = |x_2| \phi(x_2^{-1} x_1)\) if \(x_2 \neq 0\) . Show that \(\| (x_1, x_2) \|\) is an ultranorm on \(\mathbf{K}^2\) (I, p.~26, exerc. 12) and show that there does not exist any projection of norm 1 of \(\mathbf{K}^2\) on \(\mathbf{K} \times \{0\}\) .

\begin{enumerate}
\def\labelenumi{\alph{enumi})}
\setcounter{enumi}{1}
\item
  Let K be a complete non-discrete valued division ring of which the absolute value is an ultrametric and which is linearly compact. Let E be a vector space of dimension 2 on K with an ultranorm and let D be a line in E; show that for all points \(x \in \mathbf{E}\) , there exists \(y \in \mathbf{D}\) such that \(d(x, D) = d(x, y) = \| x - y \|\) (note that the intersection of D and of a ball of centre \(x\) is a ball in D).
\item
  Deduce from \(a\) ) and \(b\) ) that for a complete non-discrete valued division ring \(K\) , of which the absolute value is an ultrametric, the following properties are equivalent:
\end{enumerate}

\(\alpha)\) K is linearly compact.

β) For every ultranormed vector space E on K, for every vector subspace F of E and every continuous linear form f on F there exists a continuous linear form g on E that extends f and is such that \(\|g\| = \|f\|\) . (Reduce to the case where E is of dimension 2 and use b).) *

\exercisegroup

\begin{enumerate}
\def\labelenumi{\arabic{enumi})}
\item
  Let E be a vector space and A a convex symmetric convex subset of E. Let T, T' be two locally convex topologies on E and U, U' be the uniform structures defined by T, T' on E. In order that the uniform structure induced on A by U' should be finer than that induced by U, it is necessary and sufficient that every neighbourhood of 0 for the topology induced on A by T should be a neighbourhood of 0 for the topology induced on A by T'.
\item
  \begin{enumerate}
  \def\labelenumii{\alph{enumii})}
  \tightlist
  \item
    Give an example of a non-compact closed set in \(\mathbf{R}^2\) , whose convex envelope is not closed. b) Show that, in \(\mathbf{R}^n\) , the convex envelope of a compact set is compact (cf.~II, p.~66, exerc. 9, a)).
  \end{enumerate}
\end{enumerate}

¶ 3) Let I be the compact interval {[}0, 1{]} of R and F be the vector space \(\mathcal{C}(\mathrm{I},\mathbf{R})\) of continuous real valued functions defined in I. Let E be the product space \(\mathbf{R}^{\mathbf{F}}\) ; for all \(a\in\mathrm{I}\) , let \(\varepsilon_{a}\) be the element of E such that \(\varepsilon_{a}(f)=f(a)\) for all \(f\in\mathrm{F}\) .

\begin{enumerate}
\def\labelenumi{\alph{enumi})}
\item
  Show that, when \(x\) varies in I, the set \(\mathbf{K}\) formed by the \(\varepsilon_x\) is compact in E.
\item
  Let \(\lambda\) be an element of \(E\) such that \(\lambda(f) = \int_0^1 f(t)dt\) for all \(f\in F\) (Lebesgue measure).
\end{enumerate}

Show that, in E, \(\lambda\) belongs to the closure of the convex envelope of K but does not belong to this convex envelope (cf.~FVR, II, p.~7, prop. 5).

\begin{enumerate}
\def\labelenumi{\arabic{enumi})}
\setcounter{enumi}{3}
\tightlist
\item
  With the notations of II, p.~72, exerc. 1 suppose also that the space \(\mathbf{E}\) is locally convex.
\end{enumerate}

\begin{enumerate}
\def\labelenumi{\alph{enumi})}
\item
  There exists a positive continuous linear form \(g\) in \(\mathbf{E}\) that extends \(f\), if and only if \(f\) is bounded below in \(M \cap (\mathbf{W} + \mathbf{P})\) for at least one neighbourhood \(\mathbf{W}\) of 0 in \(\mathbf{E}\).
\end{enumerate}

\begin{enumerate}
\def\labelenumi{\alph{enumi})}
\setcounter{enumi}{1}
\tightlist
\item
  Given a point \(x \in \mathbf{E}\) , there exists a positive continuous linear form \(g\) in \(\mathbf{E}\) such that \(g(x) = 1\) , if, and only if, \(-x \notin \overline{\mathbf{P}}\) .
\end{enumerate}

\begin{enumerate}
\def\labelenumi{\arabic{enumi})}
\setcounter{enumi}{4}
\tightlist
\item
  \begin{enumerate}
  \def\labelenumii{\alph{enumii})}
  \tightlist
  \item
    Let E be an infinite dimensional normed space and T be its topology. Show that there exists on E a normed space topology \(T'\) that is strictly finer than that of T and a normed space topology \(T''\) that is strictly coarser than that of T (define the neighbourhoods of 0 for these topologies, using a basis of E put in the form \((a_{\alpha,n})\) where \(\alpha\) varies in an infinite set of indices A and n in the set of integers \(\geqslant 0\) and where \(\|a_{\alpha,n}\| = 1\) for the given norm on E).
  \end{enumerate}
\end{enumerate}

\begin{enumerate}
\def\labelenumi{\alph{enumi})}
\setcounter{enumi}{1}
\item
  Let \(p\) be the norm defining the topology \(\mathcal{T}'\) . Show that, if \(E\) is complete for the topology \(\mathcal{T}\) , then \(p\) cannot be lower semi-continuous in \(E\) for the topology \(\mathcal{T}\) (use Baire's th. cf.~III, p.~25, corollary). Deduce that the convex set \(A\) defined by the relation \(p(x) < 1\) does not contain any interior point for \(\mathcal{T}\) even though all its points are internal.
\item
  Deduce from \(b\) ), that, if \(E\) is complete for the topology \(\mathcal{T}\) , then there exists in \(E\) convex sets which are not closed for \(\mathcal{T}\) , of which the intersection with every linear variety of finite dimension is closed for \(\mathcal{T}\) (cf.~II, p.~67, exerc. 12).
\end{enumerate}

\begin{enumerate}
\def\labelenumi{\arabic{enumi})}
\setcounter{enumi}{5}
\tightlist
\item
  Let \(\mathbf{E}\) be a vector space with its finest locally convex topology.
\end{enumerate}

\begin{enumerate}
\def\labelenumi{\alph{enumi})}
\item
  Show that every vector subspace of E is closed, and that, if M, N are two subspaces that are vectorial complements in E, then E is the direct topological sum of M and N. If \((e_{t})_{t\in I}\) is a basis for E, then E is the direct topological sum of the subspaces \(Re_{t}\) .
\item
  Let \(\mathrm{F}\) be a locally convex space whose topology is also the finest locally convex topology. Show that every linear mapping of \(\mathrm{E}\) in \(\mathrm{F}\) is a strict morphism.
\end{enumerate}

\begin{enumerate}
\def\labelenumi{\arabic{enumi})}
\setcounter{enumi}{6}
\item
  \begin{enumerate}
  \def\labelenumii{\alph{enumii})}
  \tightlist
  \item
    Let A be a convex set with at least one interior point in a topological vector space E. Show that the set of internal points of A is identical with the interior of A (cf.~exerc. 5, b)). b) Show that in the normed space \(E = l^{1}(\mathbf{N})\) , the convex cone P defined in II, p.~66, exerc. 11, b), generates E but does not contain any internal point.
  \end{enumerate}
\item
  Let E be a vector space with an enumerable basis and with the finest locally convex topology. Show that, if A is a set in E whose intersection with every vector subspace of finite dimension is closed in E, then A is closed in E (cf.~exerc. 5, c)).
\item
  Let E and F be two vector spaces each with its finest locally convex topology.
\end{enumerate}

\begin{enumerate}
\def\labelenumi{\alph{enumi})}
\item
  Show that if E and F each have an enumerable basis then every bilinear mapping of \(E \times F\) in a locally convex space G is continuous (use Du Bois-Reymond's th. FVR, V, p.~53, exerc. 8).
\item
  If one of the spaces E, F has a basis with cardinal equal to that of the continuum, show that there exists a non-continuous bilinear form in \(E \times F\) . (Reduce to the case where \(E = R^{(N)}\) , \(F = R^{N}\) , so that F can be identified with \(E^{*}\) and the bilinear forms on \(E \times F\) correspond bijectively with the linear mappings of \(E^{*}\) in itself; then consider the identity mapping of \(E^{*}\) , and note that in \(R^{N}\) , a compact set for the product topology cannot be absorbent.)
\end{enumerate}

\begin{enumerate}
\def\labelenumi{\arabic{enumi})}
\setcounter{enumi}{9}
\item
  Let \((\mathrm{E}_n)\) be an infinite sequence of locally convex spaces and let E be the topological direct sum of the family \((\mathrm{E}_n)\) . Show that the topology of E is identical with the topology \(\mathcal{T}_0\) defined in I, p.~24, exerc. 14.
\item
  Let \(\mathbf{I}\) be an infinite non-enumerable set. On the vector space \(\mathrm{E} = \mathbf{R}^{(I)}\) , show that the finest locally convex topology is distinct from the topology \(\mathcal{T}_0\) defined in \(\mathbf{I}\) , p.~24, exerc. 14; for this prove that the set of the \(x = (\xi_{i})_{i\in I} \in \mathrm{E}\) such that \(|\sum_{i\in I} \xi_i| < 1\) , is open in \(\mathcal{T}\) but not in \(\mathcal{T}_0\) .
\item
  Let \(\mathbf{E}\) be a vector space with an enumerable basis \((e_n)\) . Let \(\mathbf{V}\) be the balanced convex envelope of the set of the \(e_n\) and let \(\mathbf{W}\) be the balanced convex envelope of the set of points
\end{enumerate}

\[
a _ {n} = e _ {n} + (n - 1) e _ {1} \quad (n \geqslant 1).
\]

Let \(T_{1}\) (resp. \(T_{2}\) ) be the locally convex topology on E for which a fundamental system of neighbourhoods of 0 is formed by the \(\lambda V\) (resp. \(\lambda W\) ) for \(\lambda > 0\) . Show that \(T_{1}\) and \(T_{2}\) are Hausdorff, but that the lower bound of \(T_{1}\) and \(T_{2}\) in the set of locally convex topologies on E is not Hausdorff (cf.~II, p.~80, exerc. 26).

\begin{enumerate}
\def\labelenumi{\arabic{enumi})}
\setcounter{enumi}{12}
\item
  With the hypotheses of II, §6, show that E is complete for topology T which is the inductive limit of the \(T_{n}\) , if and only if, for each integer n and every Cauchy filter F on \(E_{n}\) for the topology induced by T, there exists \(p \geqslant n\) such that F is convergent in \(E_{p}\) for the topology \(T_{p}\) .
\item
  Let E be the strict inductive limit of an increasing sequence of locally convex spaces \(E_{n}\) (II, p.~33). Show that the topology of E is the finest of the topologies compatible with the vector space structure of E, whether locally convex or not, and inducing on \(E_{n}\) a coarser topology than the given topology \(T_{n}\) . (Let \(V_{0}\) be a neighbourhood of 0 for such a topology T and \((\mathrm{V}_{n})_{n\geqslant0}\) , a sequence of neighbourhoods of 0 for T such that \(V_{n+1} + V_{n+1} \subset V_{n}\) for all \(n \geqslant 0\) ; for all \(n \geqslant 1\) , consider, in \(E_{n}\) , a convex neighbourhood \(W_{n}\) of 0 that is contained in \(E_{n} \cap V_{n}\) , and take the convex envelope of the union of the \(W_{n}\) in E.)
\item
  Let I be an infinite non-enumerable set. Let \(\mathfrak{F}(\mathrm{I})\) be the family of finite subsets of I and E the direct sum space \(\mathbf{R}^{(I)}\) . For every \(J \in \mathfrak{F}(\mathrm{I})\) , let \(F_{J}\) be the subspace \(\mathbf{R}^{\mathrm{J}}\) of E, the product of the factors whose indices belong to J, with the product topology; let \(g_{J}\) be the canonical injection of \(F_{J}\) in E. Show that there exists a topology \(\mathcal{T}_0\) on \(E\) that is compatible with the vector space structure of \(E\) , which make the \(g_{J}\) continuous and which is strictly finer than the finest locally convex topology \(\mathcal{T}\) which makes the \(g_{J}\) continuous. (Note that the set \(V\) of the \(x = (\xi_{t})_{t\in I}\) in \(E\) such that \(\sum_{i\in I}|\xi_i|^{1 / 2}\leqslant 1\) is a neighbourhood of 0 for a topology compatible with the vector
\end{enumerate}

space structure of E, not containing any absorbent, symmetric, convex set.)

\begin{enumerate}
\def\labelenumi{\arabic{enumi})}
\setcounter{enumi}{15}
\tightlist
\item
  \begin{enumerate}
  \def\labelenumii{\alph{enumii})}
  \tightlist
  \item
    Let E be a vector space with an enumerable basis. Show that the finest locally convex topology on E is the finest of the topologies on E (compatible or not with the vector space structure of E) which induces the canonical topology on every finite dimensional subspace of E.
  \end{enumerate}
\end{enumerate}

\begin{enumerate}
\def\labelenumi{\alph{enumi})}
\setcounter{enumi}{1}
\tightlist
\item
  Let \(E_{0}\) be an infinite dimensional Banach space. Let E be the vector space that is the direct sum of \(E_{0}\) of \(\mathbf{R}^{(\mathbf{N})}\) and let \(E_{p}\) be the subspace of E that is the direct sum of \(E_{0}\) and of \(R^{p}\) (identified as the product of the first p factors of \(\mathbf{R}^{(\mathbf{N})}\) ; we give to \(E_{p}\) the product topology of those of its factor, so that the topology of \(E_{p}\) is induced by that of \(E_{p+1}\) . Show that on E the inductive limit topology of those of the \(E_{p}\) is not the finest of the topologies (compatible or not with the vector-space structure of E) which induces on each \(E_{p}\) a coarser topology than that of \(E_{p}\) . We can proceed as follows :
\end{enumerate}

\(\alpha)\) Let \(q\) be a norm on \(\mathbf{E}_0\) which defines a topology strictly coarser than that of \(\mathbf{E}_0\) (II, p.~74, exerc. 5). For every \(\varepsilon > 0\) define a mapping \(f_{\varepsilon}\) of \(\mathbf{E}_0\) in \(\mathbf{R}_+\) by the relation \(f_{\varepsilon}(x) = \sup(q(x), \varepsilon - \| x \|)\) . Show that \(f_{\varepsilon}\) is continuous and \(>0\) in \(\mathbf{E}_0\) and that

\[
\inf_{\| x \| = \varepsilon} f_{\varepsilon}(x) = 0 .
\]

\(\beta)\) Let \(U\) be the subset of \(E\) formed by the \((x, (t_n))\) such that \(t_n < f_{1/n}(x)\) for all \(n\) . Show that \(U \cap E_p\) is open in \(E_p\) for all \(p\) .

\(\gamma)\) Show that if \(\mathrm{V} \subset \mathrm{U}\) is an absorbent convex set, then \(\mathrm{V} \cap \mathrm{E}_0\) cannot contain any ball with centre 0 in \(\mathrm{E}_0\) .

\begin{enumerate}
\def\labelenumi{\arabic{enumi})}
\setcounter{enumi}{16}
\tightlist
\item
  For a subset A of a commutative group G, written additively, and for each \(n > 0\) denote the set of elements of the form \(\sum_{i=1}^{n} x_i\) , where \(x_i \in \mathbf{A}\) for all \(i\) by \({}^{n}+\mathbf{A}\) . We say that the set A
\end{enumerate}

of G is convex if, for every integer n > 0 the relation \(nx \in {}^{n}+A\) implies \(x \in A\) .

\begin{enumerate}
\def\labelenumi{\alph{enumi})}
\item
  Show that if a commutative topological group G (written additively) is isomorphic to a subgroup of the additive group of a locally convex vector space (with the induced topology) then there exists a fundamental system of symmetric convex neighbourhoods of 0 in G.
\item
  Conversely, let G be a Hausdorff topological commutative group (written additively) in which there exists a fundamental system, B, of symmetric convex neighbourhoods of 0. Show that G is without torsion, and, hence, can be considered (algebraically) as a subgroup of the additive group of a vector space on the field Q (A, II, § 7.10, cor. 1 to prop. 26). For every set V ∈ B, let \(\tilde{V}\) be the set of elements rx where x ∈ V and r varies in the set of rational numbers such that \(0 \leqslant r \leqslant 1\) ; show that \(\tilde{V}\) is symmetric and convex (in the sense defined above). Deduce, further, that if there is no open subgroup of G distinct from G itself, then the sets \(\tilde{V}\) form a fundamental system of neighbourhoods of 0 for a topology compatible with the vector space structure of E on Q (Q being given its usual topology); conclude that in this case G is isomorphic to a subgroup of the additive group of a Hausdorff locally convex space.
\item
  Let G be the group \(R \times R\) ordered lexicographically (A, VI, p.~7); consider the Hausdorff topology \(\mathcal{T}_{0}(G)\) on G that is compatible with its group structure (GT, IV, § 1, exerc. 1). Show that for this topology there exists a fundamental system of symmetric convex neighbourhoods of 0, but that G is not isomorphic to any subgroup of the additive group of a Hausdorff topological vector space over R.
\end{enumerate}

\exercisegroup

\begin{enumerate}
\def\labelenumi{\arabic{enumi})}
\tightlist
\item
  \begin{enumerate}
  \def\labelenumii{\alph{enumii})}
  \tightlist
  \item
    Let E be a vector space. We say that a pointed convex cone C (of vertex 0) in E is maximal if C is a maximal element of the set of convex pointed cones of vertex 0 and \(\neq E\) , ordered by inclusion. Show that a pointed convex cone C is maximal if, and only if, it is a closed half-space defined by a hyperplane which passes through 0. To establish this result, prove successively the following properties of a maximal pointed convex cone C;
  \end{enumerate}
\end{enumerate}

α) We have C ∪ (− C) = E (argue by obtaining a contradiction).

\(\beta)\) If \(z\) is a non-internal (II, p.~26) point of C then \(-z \in C\) (same method). Deduce that C contains internal points.

\(\gamma)\) The largest vector subspace \(H = C \cap (-C)\) contained in C is a hyperplane. (Passing to the quotient space F = E/H, this reduces to demonstrating, using \(\beta)\) that if all the points of C other than the vertex are internal, then E is necessarily of dimension 1.)

\begin{enumerate}
\def\labelenumi{\alph{enumi})}
\setcounter{enumi}{1}
\tightlist
\item
  Give an example of a maximal non-pointed convex cone (in the set of non-pointed convex cones of vertex 0) which has no internal point (cf.~II, p.~65, exerc. 5).
\end{enumerate}

\begin{enumerate}
\def\labelenumi{\arabic{enumi})}
\setcounter{enumi}{1}
\item
  Let N be a hyperplane in a vector subspace M of a vector space E and let A be a convex set in E, such that all the points of \(A \cap M\) are on the same side of N and which also possesses the following property; for any \(y \neq 0\) in E, there exists \(x \in A \cap M\) such that \(x + \lambda y \in A\) for all \(\lambda\) such that \(|\lambda|\) is sufficiently small. Show that there exists then, a hyperplane H of E such that all the points of A are on the same side of H and such that \(H \cap M = N\) . (Reduce to the case \(N = \{0\}\) ; if \(a \neq 0\) belongs to \(A \cap M\) , consider the set U of pointed convex cones with vertex 0 containing A and not containing -a; show that there exists a maximal element C of U and that C is a maximal pointed convex cone (exerc. 1)). Deduce a new proof of the Hahn-Banach theorem.
\item
  Let A be a convex set in a topological vector space E and \(x_{0}\) be a point of E. Then, there exists a closed hyperplane H, containing \(x_{0}\) , and such that all the points of A lie on the same side of H if and only if there exists a non-pointed convex cone C with vertex \(x_{0}\) , which contains at least one interior point and does not meet A. (For an example of a convex set \(A \neq E\) which is not contained in any half-space defined by a hyperplane, see II, p.~65, exerc. 5.)
\item
  Let E be a normed space and A a complete convex set for the uniform structure induced by that of E.
\end{enumerate}

\begin{enumerate}
\def\labelenumi{\alph{enumi})}
\item
  Let \(x'\) be a continuous linear form on \(E\) that is bounded in \(A\) . Consider a number \(k > 0\) and the closed convex cone \(P\) in \(E\) , with vertex 0 and formed by the \(x \in E\) such that \(\| x \| \leqslant k < \langle x, x' \rangle\) ; it is pointed and proper. Show that for the order on \(E\) for which \(P\) is the set of elements \(\geqslant 0\) , the set \(A\) is inductive (use the fact that the restriction of \(x'\) to \(A\) is increasing and bounded).
\item
  Deduce from a) that the set of points of the frontier F of A which belong to a support hyperplane of A is dense in F (Bishop-Phelps th.). (For each point \(z \in F\) , consider a point \(y \in \complement A\) arbitrarily close to z, and separate y strictly from A by a closed hyperplane of equation \(\langle x, x' \rangle = \alpha\) , with \(\|x' \| = 1\) and use a) with k > 1, also exerc. 3 above.)
\end{enumerate}

\begin{enumerate}
\def\labelenumi{\arabic{enumi})}
\setcounter{enumi}{4}
\tightlist
\item
  Let \(\mathbf{A}\) be a closed convex set in \(\mathbf{R}^n\) and \(x_0\) be a point of \(\complement \mathbf{A}\) ; denote the euclidean distance in \(\mathbf{R}^n\) by \(d\) .
\end{enumerate}

\begin{enumerate}
\def\labelenumi{\alph{enumi})}
\item
  Show, without using th. 1 of II, p.~36, that there exists one and only one point \(x \in \mathbf{A}\) such that \(d(x_0, x) = d(x_0, \mathbf{A})\) , and that the hyperplane orthogonal to the line joining \(x_0\) and \(x\) , and passing through \(x\) is a support hyperplane of \(\mathbf{A}\) .
\item
  Deduce from a) a new proof of th. 1 of II, p.~36 when the space E is finite dimensional. (Reduce to the case when M is a frontier point \(x_{0}\) of A; note that the lower bound of the distance of \(x_{0}\) from support hyperplanes of A, is zero, and use the compactness of \(S_{n-1}\) .)
\end{enumerate}

\begin{enumerate}
\def\labelenumi{\arabic{enumi})}
\setcounter{enumi}{5}
\item
  Let A be a closed set in \(R^{n}\) with the following properties; for every \(x \in R^{n}\) , there exists one and only one point \(y \in A\) such that \(d(x, y) = d(x, A)\) , where d is the euclidean distance. Show that A is convex. (Argue by reductio ad absurdum, considering a closed segment with end points a, b in A containing a point \(c \in \complement A\) ; there is a closed ball B of centre c contained in \(\complement A\); consider the set B of closed balls S which contain B and whose interiors do not meet A; show that the radii of these balls is bounded above, and deduce that there exists one of these balls \(S_{0}\) whose radius \(\rho\) is the largest possible. Then get a contradiction by proving that \(S_{0}\) can only meet A in a single point, and that this implies the existence in B of a ball of radius \(> \rho\) .)
\item
  In a Hausdorff locally convex space \(\mathbf{E}\) , let \(\mathbf{A}\) be a complete, convex set and let \(\mathbf{B}\) be a precompact closed convex set such that \(\mathbf{A} \cap \mathbf{B} = \varnothing\) . Show that there exists a closed hyperplane separating A from B (argue in the completion \(\hat{\mathbf{E}}\) ). Consider the case when A is finite dimensional.
\item
  In a Hausdorff locally convex space, let A and B be two closed convex sets, without common points, such that \(C_{A} \cap C_{B} = \{0\}\) (II, p.~67, exerc. 14), and such that B is locally compact. Show that there exists a closed hyperplane separating A from B (cf.~II, p.~67, exerc. 16). Similarly, if A and B are two closed convex cones of vertex 0, such that \(A \cap B = \{0\}\) and B is locally compact, then there exists a closed hyperplane passing through 0 and separating A from B (use lemma 1 of II, p.~39).
\item
  Deduce from exerc. 8 that if V is a finite dimensional vector subspace of E and C is a closed convex cone of vertex 0 in E such that \(C \cap V = \{0\}\) , then there exists a support hyperplane of C that contains V (use lemma 1 of II, p.~39).
\item
  In the normed space \(E = l^{1}(\mathbf{N})\) of summable sequences of real numbers \(x = (\xi_{n})_{n \in \mathbf{N}}\) , let D be the line defined by the relations \(\xi_{n} = 0\) for \(n \geqslant 1\) . Show that there exist two increasing sequences \((\alpha_{n})\) , \((\beta_{n})\) of real numbers > 0 such that the convex set A defined by the inequalities \(\xi_{0} \geqslant |\alpha_{n}\xi_{n} - \beta_{n}|\) for \(n \geqslant 1\) is closed, non-bounded does not meet D and that there is no closed hyperplane separating A from D (choose \(\alpha_{n}\) and \(\beta_{n}\) so that A - D is everywhere dense).
\end{enumerate}

* 11) a) Let E be a Hilbert space and F an everywhere dense subspace of the dual E' of E that is distinct from E'; the unit ball B of E is compact for the weak topology σ(E, F), and there exists a point a of the unit sphere through which passes no closed (in σ(E, F)) support hyperplane of B.

\begin{enumerate}
\def\labelenumi{\alph{enumi})}
\setcounter{enumi}{1}
\item
  Give E the topology \(\sigma(\mathrm{E},\mathrm{F})\) and consider, in the product space \(\mathbf{G} = \mathbf{E} \times \mathbf{R}\) the set A of pairs \((x,\zeta)\) such that \(\|x\| < 1, \zeta \geqslant \|x\| / (1 - \|x\|)\) . Show that A is closed and locally compact, but that if D is the line with equation \(x = a\) in G then \(\mathrm{D} \cap \mathrm{A} = \emptyset\) and there does not exist any closed hyperplane in G separating A from D.
\item
  Show that, when we give E the topology \(\sigma(E, F)\) , there exists a continuous affine real valued function in the subspace B of E, that is not the restriction to B of a continuous affine function in E. *
\end{enumerate}

\begin{enumerate}
\def\labelenumi{\arabic{enumi})}
\setcounter{enumi}{11}
\item
  Consider in \(\mathbf{R}^3\) , the closed convex cone \(C\) defined by the relations \(\xi_1 \geqslant 0\) , \(\xi_2 \geqslant 0\) , \(\xi_3^2 \leqslant \xi_1\xi_2\) . Show that the line \(D\) of equations \(\xi_1 = 0\) , \(\xi_3 = 1\) does not meet \(C\) , but that there is no plane through the origin \(0\) containing \(D\) and not meeting \(C - \{0\}\) .
\item
  Let A and B be two closed convex sets in the space \(R^{n}\) , such that if V and W are affine linear varieties generated by A and B respectively, then no point of \(A \cap B\) is both interior to A relative to V and interior to B relative to W. Show that there exists a hyperplane separating A from B. (By taking quotients, reduce it to the case where either one of the varieties V, W is contained in the other or V and W are complementary vector subspaces in E.)
\item
  Let A be a parabolic closed convex set (II, p.~67, exerc. 17) not containing a line. Show that if B is a closed convex set not meeting A then there exists a hyperplane in \(\mathbf{R}^n\) that separates A strictly from B (if \(d\) is the Euclidean distance prove that \(d(\mathrm{A},\mathrm{B}) > 0\) ) (cf.~exerc. 12.)
\item
  Let S, T two finite sets in \(R^{n}\) with no common points, and such that \(\operatorname{Card}(S \cup T) \geqslant n + 2\) . In order that there exists a hyperplane separating S strictly from T, it is necessary and sufficient that for every finite set \(F \subset S \cup T\) of \(n + 2\) points, there exists a hyperplane separating \(F \cap S\) strictly from \(F \cap T\) (use Helly's th. (II, p.~68, exerc. 21)). Show that in this statement we cannot replace the number \(n + 2\) by \(n + 1\) , and that the statement does not extend to the case where S and T are infinite.
\item
  Let A be a compact set with interior points in \(R^{n}\) . Show that if each frontier point of A lies on at least one support hyperplane of A, then A is convex. (Obtain a contradiction, showing that if x and y are two points of A such that the segment with end points x, y is not contained in A, and if z is an interior point of A not situated on this segment, then there exists a frontier point of A, distinct from x and y in the triangle with vertices x, y, z.)
\item
  In \(R^{n}\) , let A be a symmetric convex set of which 0 is an interior point and of which the frontier does not contain any genuine segment. Let H be a homogeneous hyperplane and D a line complementary to H. Show that there exists a point \(a \in H \cap A\) such that at a there is a hyperplane of support to A that is parallel to D.
\item
  In a topological vector space \(\mathbf{E}\) , let \(\mathrm{A}_i\) ( \(1 \leqslant i \leqslant n\) ) be \(n\) open non-empty convex sets. a) Show that if the union of the \(\mathrm{A}_i\) is distinct from \(\mathbf{E}\) , then every point \(x \in \mathbf{E}\) not belonging to any of the \(\mathrm{A}_i\) , belongs to a closed linear variety of codimension \(n\) , that contains \(x\) and does not meet any of the \(\mathrm{A}_i\) (argue by induction on \(n\) ).
\end{enumerate}

\begin{enumerate}
\def\labelenumi{\alph{enumi})}
\setcounter{enumi}{1}
\tightlist
\item
  If the intersection of the \(\mathbf{A}_i\) is empty, show that there exists, in E, a closed linear variety of codimension \(n - 1\) that does not meet any of the \(\mathbf{A}_i\) (same method).
\end{enumerate}

\begin{enumerate}
\def\labelenumi{\arabic{enumi})}
\setcounter{enumi}{18}
\tightlist
\item
  Let C, C' be two closed convex sets in a Hausdorff topological vector space E that are strictly separated by a closed hyperplane H. Let H' be a closed hyperplane of support to both C and C' such that C and C' lie on the same side of H'. Show that H' is the only hyperplane with these properties which contains H ∩ H' and that H ∩ H' is a support hyperplane of the trace P on H of the convex envelope of C ∪ C'. Conversely, if C and C' are compact, then for every support hyperplane D of P in H, there exists a hyperplane H' that supports both C and C', which contains D and such that C and C' lie on the same side of H'.
\end{enumerate}

¶ 20) Let A, B be two disjoint closed convex sets in a Hausdorff locally convex space E and let H be a closed hyperplane separating A from B; suppose that \(A \cap H \neq \varnothing\) and that the intersection of \(A \cap H\) and of every line is compact. Show that, if A or B is locally compact, then there exists a neighbourhood V of 0 in E such that \((A + V) \cap B\) is empty. (Consider two cases according to whether A or B is locally compact; in the first case, note that there exists a hyperplane \(H'\) parallel to H such that, if S is the set of points between H and \(H'\) , then \(A \cap S\) is compact. In the second case, suppose for example that \(0 \in B \cap H\) ; for every neighbourhood V of 0 in E, consider the set \((A + V) \cap B\) and consider successively the case where this set is relatively compact for at least one V or the case when this is not so, as in exerc. 16 of II, p.~67.)

¶ 21) a) In \(\mathbf{R}^n\) let \(a_i (1 \leqslant i \leqslant n + 1)\) be \(n + 1\) points that are affinely independent. Denote the convex envelope of the \(a_i\) by \(\mathbf{S}\) and the convex envelope of the \(a_i\) with \(i \neq k\) by \(\mathrm{F}_k\) for \(1 \leqslant k \leqslant n + 1\) . For each \(k\) let \(\mathrm{C}_k\) be a compact convex set containing \(\mathrm{F}_k\) , and suppose that

S is contained in the union of the \(C_k\) ; show then that \(\bigcap_{k=1}^{n+1} C_k \neq \emptyset\) . (Argue by reductio ad

absurdum and induction on \(n\) , considering the intersection \(\mathbf{C}_{n+1}'\) of the \(\mathbf{C}_i\) with indices \(i \leqslant n\) and supposing that \(\mathbf{C}_{n+1} \cap \mathbf{C}_{n+1}' = \emptyset\) , which would allow the strict separation of the two convex sets by a hyperplane.)

\begin{enumerate}
\def\labelenumi{\alph{enumi})}
\setcounter{enumi}{1}
\tightlist
\item
  Let X be a compact convex set in a Hausdorff topological vector space E, and \((\mathrm{C}_{\lambda})_{\lambda\in\mathrm{L}}\) a family of compact convex sets contained in X, such that for every set H ⊂ L having n (resp. m) elements, the intersection (resp. the union) of the \(C_{\lambda}\) with indices \(\lambda\in H\) is not empty (resp. is equal to X). Show that if \(m\leqslant n+1\) the intersection \(\bigcap_{\lambda\in L}C_{\lambda}\) is not empty. (This is effectively
\end{enumerate}

proving that for any finite set H of \(p \geqslant m\) indices of L, we have \(\bigcap_{\lambda \in H} C_{\lambda} \neq \varnothing\) . Argue by induc-

tion on \(p\) assuming that the result has been proved for \(p - 1\) indices. Argue then by reductio ad absurdum, considering for each index \(i \in \mathbf{H}\) a point \(a_i \in \bigcap_{\lambda \in \mathbf{H} - \{i\}} \mathrm{C}_\lambda\) , and showing by the

aid of Helly's th. (II, p.~68, exerc. 21) that the \(a_{i}\) generate a linear variety of dimension \(p - 1\) , then finally apply \(a\) .)

¶ 22) In \(\mathbf{R}^n\) let \((\mathbf{C}_i)_{1 \leqslant i \leqslant m}\) be a finite family of closed convex cones with vertex 0, such that the sum of any \(n\) of them is distinct from \(\mathbf{R}^n\) . Show that there exists a hyperplane H, passing through 0, such that, for any index \(i\) no pair of points of \(\mathbf{C}_i\) are strictly separated by H. (Distinguish two cases:

\(\alpha)\) Either there exists a number \(r < n\) and \(r\) indices, say, 1, 2, \ldots, \(r\) such that \(C_i\) for \(i \leqslant r\) generates a cone which contains a vector subspace V of dimension \(\geqslant r\) . Argue by induction on \(n\) , projecting on the orthogonal to V.

β) Or, for all r < n, any r of the \(C_{i}\) generate a cone C such that the maximal vector subspace \(C \cap (-C)\) contained in C is of dimension < r. Consider then a set of the cones \(C_{i}\) , maximal with respect to being contained in a half-space; there are at least n cones in a maximal set. Let \(\Gamma\) be the cone generated by the union of the cones belonging to this maximal set. If \(C_{j}\) is a cone which does not belong to the maximal set considered, show that \(C_{i} \subset -\Gamma\) . For this, argue by reductio ad absurdum, showing that in the contrary case there exists a frontier point of \(-\Gamma\) (relative to the vector subspace generated by \(\Gamma\) ) that is interior to \(C_{j}\) (relative to the vector subspace generated by \(C_{j}\) ). Write such a point as the sum of the least number s of vectors, of which each belongs to a cone \(-C_{i}\) , among those \(C_{i}\) used in defining \(\Gamma\) ; then \(s \leqslant n - 1\) . Prove finally that these cones and \(C_{j}\) generate a convex cone containing an \(s + 1\) dimensional vector subspace, contradicting the hypothesis; for this use exerc. 4 of II, p.~65.)

\begin{enumerate}
\def\labelenumi{\arabic{enumi})}
\setcounter{enumi}{22}
\item
  Let \(\mathbf{E}\) be a topological vector space, and let \(\mathcal{T}\) be the locally convex topology on \(\mathbf{E}\) that is the finest of all those that are coarser than the given topology \(\mathcal{T}_0\) on \(\mathbf{E}\) . If \(\mathbf{F}\) is a locally convex space, then the continuous linear mappings of \(\mathbf{E}\) in \(\mathbf{F}\) are the same for \(\mathcal{T}_0\) as for \(\mathcal{T}\) . There exists a continuous linear form on \(\mathbf{E}\) that is distinct from the null form if, and only if, there exists a neighbourhood of 0 for \(\mathcal{T}_0\) whose convex envelope is not everywhere dense (for \(\mathcal{T}_0\) ) (cf.~I, p.~25, exerc. 4).
\item
  Let \(\mathbf{E}\) be an infinite dimensional, metrisable, locally convex space.
\end{enumerate}

\begin{enumerate}
\def\labelenumi{\alph{enumi})}
\item
  Show that there exists a sequence \((a_{n})\) of points of E tending to 0 and a decreasing sequence \((\mathbf{L}_n)\) , of closed vector subspaces of E, such that \(\mathbf{L}_n\) is of codimension \(n\) in E and that, for all \(n\) , the point \(a_{n}\) belongs to \(\mathbf{L}_n - \mathbf{L}_{n + 1}\) .
\item
  Suppose further that E is complete. Show that we can then find sequences \((a_{n})\) and \((\mathrm{L}_{n})\) verifying the conditions a), and such that in addition, for every bounded sequence of real number \((\lambda_{n})\) , the series, whose general term is \(\lambda_{n}a_{n}\) , is commutatively convergent in E, and that the linear mapping \((\xi_{n})\mapsto\sum_{n}\xi_{n}a_{n}\) of the Banach space \(\mathcal{B}(\mathbf{N})\) in E is injective and continuous.
\item
  Deduce from b) that when E is an infinite dimensional Fréchet space then every basis of E on R has cardinal at least equal to \(2^{\mathrm{Card}(\mathbf{N})}\) (cf.~I, p.~22, exerc. 5).
\end{enumerate}

If there exists an enumerable set that is dense in \(\mathbf{E}\) , then every basis of \(\mathbf{E}\) has the cardinal of the continuum.

\begin{enumerate}
\def\labelenumi{\arabic{enumi})}
\setcounter{enumi}{24}
\item
  Let E be an infinite dimensional Fréchet space of enumerable type (therefore having an enumerable everywhere dense subset) (cf.~I, p.~25, exerc. 1). Show that there exists an everywhere dense hyperplane H of E which meets every closed, infinite dimensional linear variety of E. (Use the existence of a basis, having the cardinal of the continuum, in each of the direction subspaces of these varieties (exerc. 24, c)) and the fact that the set of closed, infinite dimensional, linear varieties of E also has the cardinal of the continuum (GT, IX, § 5, exerc. 17); then apply a method of construction of a linear form on E following from S, III, § 6, exerc. 24.) The hyperplane H does not contain any infinite dimensional, closed, vector subspace.
\item
  Let \(\mathbf{E}\) be an infinite dimensional Fréchet space of enumerable type.
\end{enumerate}

\begin{enumerate}
\def\labelenumi{\alph{enumi})}
\item
  Show that there exists a sequence \((a_{n})\) of linearly independent elements of E such that each sequence \((a_{2n})\) and \((a_{2n + 1})\) is total (use exerc. 24, c)).
\item
  Let \(\mathbf{F}\) be the vector subspace of \(\mathbf{E}\) generated by the \(a_{2n + 1}\) ( \(n \in \mathbf{N}\) ). For every \(n > 0\) let \(\mathbf{M}_n\) be the subspace generated by the \(a_{2k}\) with \(k \leqslant n\) . For each \(n\) , let \(\phi_n\) be the restriction to \(\mathbf{F}\) of the canonical homomorphism of \(\mathbf{E}\) on \(\mathrm{E / M}_n\) , and let \(\mathcal{T}_n\) be the topology on \(\mathbf{F}\) which is the inverse image under \(\phi_n\) of the quotient topology, on \(\mathrm{E / M}_n\) . Show that each of the topologies \(\mathcal{T}_n\) on \(\mathbf{F}\) is a Hausdorff locally convex topology, but that the lower bound of the \(\mathcal{T}_n\) in the set of locally convex topologies on \(\mathbf{F}\) is the coarsest topology on \(\mathbf{F}\) .
\end{enumerate}

* c) Take E to be a Hilbert space; show that we can choose the sequence \((a_{n})\) so that if G is the closed vector subspace generated by the \(a_{4n + 1}\) , then G has infinite codimension and so that the images of the \(a_{2n}\) and the \(a_{4n + 3}\) in E/G are still linearly independent. Write \(\mathbf{G}_n\) for the subspace of E that is the sum of G and of the subspace generated by the \(a_{4k + 3}\) with \(k\leqslant n\) , and give to \(G_{n}\) the topology which is the inverse image under the canonical mapping restricted to \(G_{n}\) , of the quotient topology on \(E/M_{n}\) . Show that the sequence \((\mathbf{G}_{n})\) is an inductive system of topological vector spaces such that \(G_{n}\) is closed in \(G_{n+1}\) for the topology of \(G_{n+1}\) , but that \(G_{n}\) is not closed in the inductive limit space of this sequence. *

\begin{enumerate}
\def\labelenumi{\arabic{enumi})}
\setcounter{enumi}{26}
\tightlist
\item
  Let E, F be two Hausdorff topological vector spaces, and X (resp. Y) a compact convex set in E (resp. F). Let f be a real valued function defined in \(X \times Y\) with the following properties:
\end{enumerate}

\begin{enumerate}
\def\labelenumi{(\roman{enumi})}
\item
  For all \(x \in \mathbf{X}\) , the mapping \(y \mapsto f(x, y)\) is lower semi-continuous in \(\mathbf{Y}\) , and for all \(c \in \mathbf{R}\) , the set of the \(y \in \mathbf{Y}\) such that \(f(x, y) \leqslant c\) is convex.
\item
  For all \(y \in \mathbf{Y}\) , the mapping \(x \mapsto f(x, y)\) is upper semi-continuous in \(\mathbf{X}\) , and for all \(c \in \mathbf{R}\) , the set of the \(x \in \mathbf{X}\) such that \(f(x, y) \geqslant c\) is convex.
\end{enumerate}

Show that, in these conditions, we have

\[
\sup _ {x \in X} (\inf _ {y \in Y} f (x, y)) = \inf _ {y \in Y} (\sup _ {x \in X} f (x, y)).
\]

(Argue by reductio ad absurdum, supposing that there exists a number \(c\) such that

\[
\sup _ {x \in \mathbf {X}} (\inf _ {y \in \mathbf {Y}} f (x, y)) <   c <   \inf _ {y \in \mathbf {Y}} (\sup _ {x \in \mathbf {X}} f (x, y)).
\]

For all \(x \in X\) (resp. all \(y \in Y\) ) let \(A_{x}\) be the set of \(y \in Y\) such that \(f(x, y) > c\) (resp. \(B_{y}\) the set of the \(x \in X\) such that \(f(x, y) < c\) ), which is open in Y (resp. in X); the \(A_{x}\) (resp. the \(B_{y}\) ) form a covering of Y (resp. X) when x varies in X (resp. y varies in Y). Show that there exist two finite sets \(X_{0} \subset X\) , \(Y_{0} \subset Y\) such that: \(1^{\circ}\) for all y belonging to the convex envelope \(B_{0}\) of \(Y_{0}\) , there exists \(x \in X_{0}\) such that \(f(x, y) > c\) , and \(X_{0}\) is minimal for this property; \(2^{\circ}\) for all x belonging to the convex envelope \(A_{0}\) of \(X_{0}\) , there exists \(y \in Y_{0}\) such that \(f(x, y) < c\) , and \(Y_{0}\) is minimal for this property. Then for all \(y \in Y_{0}\) , let \(C_{y}\) be the set of \(x \in A_{0}\) such that \(f(x, y) \geqslant c\) ; using exerc. 21, a) of II, p.~79, show that the intersection of the \(C_{y}\) for \(y \in Y_{0}\) is not empty. Proceed in the same way in \(B_{0}\) and obtain a contradiction.)

¶ 28) Let X be a compact convex subset of E a Hausdorff locally convex space, and let f be an upper semi-continuous convex function in X. Show that the set L of continuous convex functions g in X such that \(g(x) > f(x)\) for all \(x \in X\) is decreasing directed and that its lower envelope is equal to f.~(Let u, v be elements of L. To construct an element of L which is less than u and v, use reasoning analogous to that of prop. 6, II, p.~40. Interpret the set \(K_{1}\) analogous to the set K in this argument as the set of points situated above the graph of a lower semi continuous function that is less than u and v and strictly larger than f at every point; apply prop. 5 of II, p.~39 and Dini's th. to this function. To show that the lower envelope of L is f, note that f is bounded above by a constant b; \((x, t)\) being a point of E × R situated above the graph of f, let \(K'\) be the convex envelope of \(\{(x, t)\} \cup (X' \times \{b\})\) , where \(X'\) is a convenient compact neighbourhood of x in X; argue with \(K'\) as above for \(K_{1}\) .)

\begin{enumerate}
\def\labelenumi{\arabic{enumi})}
\setcounter{enumi}{28}
\item
  Let X be a compact convex set in a Hausdorff locally convex space E. Let u be a lower semi-continuous convex function in X and v an upper semi-continuous concave function in X such that \(u(x) > v(x)\) for all \(x \in X\) . Then there exists an affine linear function f that is continuous in E and such that \(v(x) < f(x) < u(x)\) for all \(x \in X\) .
\item
  Let \(\mathbf{X}\) be a compact convex set of a Hausdorff locally convex space \(\mathbf{E}\) . Show that the set of lower semi-continuous convex functions in \(\mathbf{X}\) is a lattice.
\end{enumerate}
